\textbf{Phần 1. Trắc nghiệm lựa chọn}
\setcounterpageref{ex}{0}

\begin{ex} %[3D2V27]
Giảm 48 đi 4 lần ta được kết quả là
\choice
{$16$}
{$14$}
{\True $12$}
{$44$}
\loigiai{
Giảm 48 đi 4 lần: $48 : 4 = 12$.
\textbf{Đáp án của câu hỏi là C}
}
\end{ex}

\begin{ex} %[3D2V27]
Muốn giảm một số đi một số lần, ta làm như sau:
\choice
{\True Lấy số đó chia cho số lần}
{Lấy số đó trừ đi số lần}
{Lấy số đó nhân với số lần}
{Lấy số đó cộng với số lần}
\loigiai{
Quy tắc toán học: Muốn giảm một số đi nhiều lần, ta lấy số đó chia cho số lần.
\textbf{Đáp án của câu hỏi là A}
}
\end{ex}

\begin{ex} %[3D2V27]
Một cuộn dây dài $72\text{ m}$, sau khi sử dụng thì độ dài cuộn dây giảm đi 3 lần. Cuộn dây còn lại dài bao nhiêu mét?
\choice
{$18\text{ m}$}
{$21\text{ m}$}
{$36\text{ m}$}
{\True $24\text{ m}$}
\loigiai{
Độ dài còn lại là: $72 : 3 = 24\text{ (m)}$.
\textbf{Đáp án của câu hỏi là D}
}
\end{ex}

\begin{ex} %[3D2V27]
Giảm số 60 đi 5 lần rồi cộng thêm 28 ta được số nào?
\choice
{$36$}
{\True $40$}
{$42$}
{$38$}
\loigiai{
$60 : 5 = 12$; $12 + 28 = 40$.
\textbf{Đáp án của câu hỏi là B}
}
\end{ex}

\begin{ex} %[3D2V27]
Một bao gạo nặng $54\text{ kg}$. Sau khi bán một số ki-lô-gam gạo thì số gạo trong bao giảm đi 6 lần. Người ta đã bán đi bao nhiêu ki-lô-gam gạo?
\choice
{$9\text{ kg}$}
{$48\text{ kg}$}
{\True $45\text{ kg}$}
{$36\text{ kg}$}
\loigiai{
Số gạo còn lại: $54 : 6 = 9\text{ kg}$. Số gạo đã bán: $54 - 9 = 45\text{ (kg)}$.
\textbf{Đáp án của câu hỏi là C}
}
\end{ex}

\begin{ex} %[3D2V27]
Số thứ nhất là 84. Số thứ hai giảm đi 4 lần so với số thứ nhất. Tổng của hai số là
\choice
{\True $105$}
{$95$}
{$100$}
{$112$}
\loigiai{
Số thứ hai là: $84 : 4 = 21$. Tổng hai số là: $84 + 21 = 105$.
\textbf{Đáp án của câu hỏi là A}
}
\end{ex}

\begin{ex} %[3D2V27]
Một số giảm đi 3 lần rồi bớt đi 6 thì được 14. Số ban đầu là
\choice
{$54$}
{$48$}
{$66$}
{\True $60$}
\loigiai{
Số đó giảm 3 lần là: $14 + 6 = 20$. Số ban đầu là: $20 \times 3 = 60$.
\textbf{Đáp án của câu hỏi là D}
}
\end{ex}

\begin{ex} %[3D2V27]
Khẳng định nào sau đây là ĐÚNG khi phân biệt giữa 'giảm đi 4 lần' và 'giảm đi 4 đơn vị' đối với số 36?
\choice
{Giảm 4 lần được 32, giảm 4 đơn vị được 9}
{\True Giảm 4 lần được 9, giảm 4 đơn vị được 32}
{Cả hai đều cho kết quả là 9}
{Cả hai đều cho kết quả là 32}
\loigiai{
Giảm 4 lần là phép chia: $36 : 4 = 9$. Giảm 4 đơn vị là phép trừ: $36 - 4 = 32$.
\textbf{Đáp án của câu hỏi là B}
}
\end{ex}

\begin{ex} %[3D2V27]
Tuổi của bà là 72 tuổi. Tuổi của mẹ giảm đi 2 lần so với tuổi bà, tuổi của cháu giảm đi 6 lần so với tuổi mẹ. Tuổi của cháu là
\choice
{$8$ tuổi}
{$12$ tuổi}
{\True $6$ tuổi}
{$7$ tuổi}
\loigiai{
Tuổi của mẹ: $72 : 2 = 36$ tuổi. Tuổi của cháu: $36 : 6 = 6$ tuổi.
\textbf{Đáp án của câu hỏi là C}
}
\end{ex}

\begin{ex} %[3D2V27]
Một thùng nước có $64\ell$. Sau khi tháo bớt nước ra thì lượng nước trong thùng giảm đi 4 lần. Lượng nước đã tháo ra ngoài là
\choice
{\True $48\ell$}
{$16\ell$}
{$32\ell$}
{$52\ell$}
\loigiai{
Lượng nước còn lại: $64 : 4 = 16\text{ (\ell)}$. Lượng nước đã tháo ra: $64 - 16 = 48\text{ (\ell)}$.
\textbf{Đáp án của câu hỏi là A}
}
\end{ex}

\begin{ex} %[3D2V27]
Số chẵn lớn nhất có hai chữ số giảm đi 3 lần thì được kết quả là
\choice
{$33$}
{$32$}
{$31$ (dư $2$)}
{\True $32$ (dư $2$)}
\loigiai{
Số chẵn lớn nhất có hai chữ số là 98. Ta có: $98 : 3 = 32$ (dư 2).
\textbf{Đáp án của câu hỏi là D}
}
\end{ex}

\begin{ex} %[3D2V27]
Nếu giảm cạnh của một hình vuông đi 3 lần thì chu vi của hình vuông đó
\choice
{Giảm đi 4 lần}
{\True Giảm đi 3 lần}
{Giảm đi 12 lần}
{Không thay đổi}
\loigiai{
Chu vi hình vuông bằng cạnh nhân 4. Cạnh giảm đi 3 lần thì chu vi cũng giảm đi 3 lần.
\textbf{Đáp án của câu hỏi là B}
}
\end{ex}

\textbf{Phần 2. Trắc nghiệm đúng sai}
\setcounterpageref{ex}{0}

\begin{ex} %[3D2V27]
Xét tính đúng sai của các nhận định về việc giảm một số đi nhiều lần:
\choiceTF
{\True Muốn giảm số $x$ đi $n$ lần ta tính $x : n$.}
{Giảm số 40 đi 5 đơn vị cũng giống như giảm số 40 đi 5 lần.}
{\True Nếu số ban đầu chia hết cho số lần thì kết quả nhận được là một số tự nhiên.}
{Giảm 90 đi 3 lần rồi giảm tiếp 3 lần nữa thì tương đương với giảm 90 đi 6 lần.}
\loigiai{
\begin{itemize}
    \item Ý a) Đúng: Định nghĩa.
    \item Ý b) Sai: Sai vì giảm 5 đơn vị là $40 - 5 = 35$, giảm 5 lần là $40 : 5 = 8$.
    \item Ý c) Đúng: Định lí chia hết.
    \item Ý d) Sai: Sai vì giảm 3 lần rồi 3 lần nữa là chia cho 9, không phải chia cho 6.
\end{itemize}
\begin{itemize}
    \item a) \textbf{Đúng}.
    \item b) \textbf{Sai}.
    \item c) \textbf{Đúng}.
    \item d) \textbf{Sai}.
\end{itemize}
}
\end{ex}

\begin{ex} %[3D2V27]
Một người nuôi 48 con thỏ. Sau khi bán đi một số thỏ thì số thỏ còn lại giảm đi 4 lần so với ban đầu:
\choiceTF
{\True Số thỏ còn lại là 12 con.}
{\True Số thỏ đã bán là 36 con.}
{Số thỏ đã bán ít hơn số thỏ còn lại.}
{\True Số thỏ đã bán gấp 3 lần số thỏ còn lại.}
\loigiai{
\begin{itemize}
    \item Ý a) Đúng: $48 : 4 = 12$ con.
    \item Ý b) Đúng: $48 - 12 = 36$ con.
    \item Ý c) Sai: Sai vì 36 con nhiều hơn 12 con.
    \item Ý d) Đúng: $36 : 12 = 3$ lần.
\end{itemize}
\begin{itemize}
    \item a) \textbf{Đúng}.
    \item b) \textbf{Đúng}.
    \item c) \textbf{Sai}.
    \item d) \textbf{Đúng}.
\end{itemize}
}
\end{ex}

\textbf{Phần 3. Trả lời ngắn}
\setcounterpageref{ex}{0}

\begin{ex} %[3D2V27]
Giảm số 96 đi 6 lần ta được số nào?
\shortans{16}
\loigiai{
\textbf{Đáp số:} $16$\\
$96 : 6 = 16$.
}
\end{ex}

\begin{ex} %[3D2V27]
Một tấm vải dài $84\text{ m}$. Sau khi may quần áo, độ dài tấm vải giảm đi 4 lần. Hỏi tấm vải còn lại bao nhiêu mét?
\shortans{21}
\loigiai{
\textbf{Đáp số:} $21$\\
Tấm vải còn lại: $84 : 4 = 21\text{ (m)}$.
}
\end{ex}

\begin{ex} %[3D2V27]
Một số giảm đi 5 lần thì được 17. Tìm số đó.
\shortans{85}
\loigiai{
\textbf{Đáp số:} $85$\\
Số đó là: $17 \times 5 = 85$.
}
\end{ex}

\begin{ex} %[3D2V27]
Thùng dầu có $72\ell$. Sau một ngày sử dụng lượng dầu trong thùng giảm đi 3 lần. Hỏi người ta đã sử dụng bao nhiêu lít dầu?
\shortans{48}
\loigiai{
\textbf{Đáp số:} $48$\\
Dầu còn lại: $72 : 3 = 24\ell$. Dầu đã dùng: $72 - 24 = 48\text{ (\ell)}$.
}
\end{ex}

\begin{ex} %[3D2V27]
Giảm số 80 đi 4 lần rồi thêm vào 35 đơn vị ta được kết quả bằng bao nhiêu?
\shortans{55}
\loigiai{
\textbf{Đáp số:} $55$\\
$80 : 4 = 20$; $20 + 35 = 55$.
}
\end{ex}

\begin{ex} %[3D2V27]
Đoạn thẳng $AB = 60\text{ cm}$. Đoạn thẳng $CD$ giảm đi 5 lần so với $AB$. Đoạn thẳng $MN$ dài gấp đôi $CD$. Tính độ dài đoạn thẳng $MN$ theo xăng-ti-mét.
\shortans{24}
\loigiai{
\textbf{Đáp số:} $24$\\
$CD = 60 : 5 = 12\text{ cm}$. $MN = 12 \times 2 = 24\text{ (cm)}$.
}
\end{ex}
